\section{Parallel Execution and Element-local Design}\label{sec:design overview}

This section describes the partitioning and parallelization strategy shared by all compressors evaluated in this work and motivates the element-local design of our two new compressors, FDWZ and MRCZ.

\subsection{Partitioned Parallel Execution}

All evaluated compressors are executed through a common partitioned processing layer. Before invoking a compressor, this layer divides each field into contiguous partitions of elements. An OpenMP thread processes each partition by invoking FDWZ or MRCZ, or by calling a baseline compressor through its API. The returned partition bitstreams are then concatenated to form the compressed field.

\subsection{Element-local Design of FDWZ and MRCZ}

FDWZ and MRCZ both adopt an element-local design for their main modeling steps. This choice is motivated by the following considerations.

\begin{enumerate}
\item \textit{Cross-element Prediction}: Prediction beyond an element can follow either the physical mesh connectivity or the stored element order. Sections~\ref{sec:FDWZ:cross-interface} and~\ref{sec:FDWZ:element-index-dimension} show that neither alternative to element-local interpolation provides a consistent prediction advantage for FDWZ.
\item \textit{Irregular Connectivity}: Accessing physically adjacent elements requires identifying their shared faces and aligning their local coordinate systems. This process is complex and time-consuming.
\item \textit{Transform Locality}: The transform used by MRCZ operates naturally on the tensor-product arrangement within each element. This arrangement is not preserved across element boundaries, making an element-local transform the natural design.
\end{enumerate}
